% Source printed pp. 59-65.
\section{Hilb 与 Quot 的射影性}\label{src:2.4}
\begin{epigraph}
射影几何就是全部几何.
\epigraphauthor{Arthur Cayley}
\end{epigraph}
我们证明 Hilbert 概形和 Quot 概形的射影性(\bookref{2.1.2}~与~\bookref{2.1.3}).
\subsection{线丛直和 \texorpdfstring{$F$}{F} 的特殊情形}\label{src:2.4.1-subsection}
设 $S$ 为 Noether 概形, $\pi:\mathbb P_S^n\to S$ 为相对射影空间.
先考虑凝聚层 $F:=\mathcal O_{\mathbb P_S^n}(-l)^{\oplus r}$ 的特殊情形, 其中 $l$ 和 $r\geq1$ 为整数.
对每个整数 $d\geq0$, 可以尝试定义函子态射
\[
\begin{aligned}
\Quot^P(F/\mathbb P_S^n)&\longrightarrow\Gr(P(d),\pi_*F(d)),\\
[F_T\twoheadrightarrow Q]&\longmapsto
[\pi_{T,*}F_T(d)\to\pi_{T,*}Q(d)],
\end{aligned}
\]
其中 $[F_T\twoheadrightarrow Q]$ 为 $S$-概形 $T$ 上的对象.
\begin{enumerate}[步骤 1: ]
\item\label{step:quot-1}
\textit{存在整数 $m_0$, 使对 $d\geq m_0$, 
$\Quot^P(F/\mathbb P_S^n)\to\Gr(P(d),\pi_*F(d))$ 良定义.}
\begin{proof}
由于 $F=\mathcal O_{\mathbb P_S^n}(-l)^{\oplus r}$, 对 $d\geq l$, 
$\pi_*F(d)=(\operatorname{Sym}^{d-l}\Gamma(S,\mathcal O_S))^{\oplus r}$
为向量丛, 且其构造与沿 $T\to S$ 的基变换交换; 对所有 $i>0$, $R^i\pi_*F(d)=0$.
对每个域值点 $s\in S(k)$ 和 Hilbert 多项式为 $P$ 的商
$F_s=\mathcal O_{\mathbb P_k^n}(-l)^{\oplus r}\twoheadrightarrow Q$, 
令 $K_s=\ker(F_s\to Q)$.
则 $K(l)$ 为 $\mathcal O_{\mathbb P_k^n}^{\oplus r}$ 的子层, 其 Hilbert 多项式由 $P$ 和 $\mathcal O_{\mathbb P_k^n}^{\oplus r}$ 的 Hilbert 多项式确定.
因此可由正则性的有界性(\bookref{2.3.8})选取 $m_0\geq l$, 使每个这样的核 $K_s$ 均为 $m_0$-正则.
必要时增大 $m_0$, 还可使 $F_s$ 为 $m_0$-正则.
习题~\bookref{2.3.4}~表明每个商 $Q$ 均为 $m_0$-正则.
由族中的正则性(\bookref{2.3.18}(1)), 对每个 $d\geq m_0$ 及在 $T$ 上平坦, Hilbert 多项式为 $P$ 的商层 $F_T\twoheadrightarrow Q$, 
$\pi_{T,*}Q(d)$ 为 $T$ 上秩为 $P(d)$ 的向量丛, 且其构造与基变换交换.
由于 $F_T$ 和 $Q$ 在 $T$ 上平坦, $K:=\ker(F_T\to Q)$ 也在 $T$ 上平坦.
每个纤维 $K_s$ 为 $m_0$-正则, 故族中的正则性(\bookref{2.3.18}(2))表明: 对 $d\geq m_0$, $R^1\pi_{T,*}K(d)=0$.
因此 $\pi_{T,*}F_T(d)\to\pi_{T,*}Q(d)$ 满射, 且~\bookeqref{2.4.1}~良定义.
\end{proof}
\item\label{step:quot-2}
\textit{存在整数 $m_0$, 使对 $d\geq m_0$, 
$\Quot^P(F/\mathbb P_S^n)\to\Gr(P(d),\pi_*F(d))$ 为单态射.}
\begin{proof}
设 $F_T\twoheadrightarrow Q$ 为在概形 $T$ 上平坦, Hilbert 多项式为 $P$ 的商层.
核 $K=\ker(F_T\to Q)$ 也在 $T$ 上平坦.
按步骤~\ref{step:quot-1}~选取 $m_0$, 则对每个 $s\in S$, 纤维 $K_s$, $(F_T)_s$ 和 $Q_s$ 均为 $m_0$-正则.
由族中的正则性(\bookref{2.3.18}(2)), 对 $d\geq m_0$, $R^1\pi_{T,*}K(d)=0$, 得到短正合列的交换图
\[
\begin{tikzcd}[column sep=small,row sep=large,font=\small]
0\arrow[r]&\pi_T^*\pi_{T,*}K(d)\arrow[r]\arrow[d]
\arrow[dr,"\alpha"]&
\pi_T^*\pi_{T,*}F_T(d)\arrow[r]\arrow[d]&
\pi_T^*\pi_{T,*}Q(d)\arrow[r]\arrow[d]&0\\
0\arrow[r]&K(d)\arrow[r]&F_T(d)\arrow[r]&Q(d)\arrow[r]&0.
\end{tikzcd}
\]
由族中的正则性(\bookref{2.3.18}(3)), 每个竖直映射均满射.
所以 $Q(d)=\operatorname{coker}(\alpha)$ 由 $\pi_{T,*}F_T(d)\to\pi_{T,*}Q(d)$ 唯一确定.
\end{proof}
\item\label{step:quot-3}
\textit{存在整数 $m_0$, 使对 $d\geq m_0$, 
$\Quot^P(F/\mathbb P_S^n)\to\Gr(P(d),\pi_*F(d))$ 为局部闭浸入.}
\begin{proof}
设 $T\to\Gr(P(d),\pi_*F(d))$ 为由向量丛商
$(\pi_*F(d))_T=\pi_{T,*}F_T(d)\twoheadrightarrow V$
确定的映射, 记其核为 $K$.
考虑纤维积
\[
\begin{tikzcd}[column sep=large]
Y\arrow[r]\arrow[d]\arrow[dr,phantom,"\square"]&
T\arrow[d,"{[\pi_{T,*}F_T(d)\twoheadrightarrow V]}"]\\
\Quot^P(F/\mathbb P_S^n)\arrow[r]&\Gr(P(d),\pi_*F(d)).
\end{tikzcd}
\]
在 $\mathbb P_T^n$ 上, 有层的交换图
\[
\begin{tikzcd}[column sep=small,row sep=large]
0\arrow[r]&\pi_T^*K\arrow[r]\arrow[dr,"\alpha"]&
\pi_T^*\pi_{T,*}F_T(d)\arrow[r]\arrow[d]&
\pi_T^*V\arrow[r]\arrow[d]&0\\
&&F_T(d)\arrow[r]&Q(d)\arrow[r]&0,
\end{tikzcd}
\]
其中 $Q$ 定义为使 $Q(d)=\operatorname{coker}(\alpha)$.
由于 $K$ 为 $T$ 上的向量丛, $K=\pi_{T,*}\pi_T^*K$, 故 $\pi_{T,*}\alpha$ 等同于包含映射 $K\hookrightarrow\pi_{T,*}F_T(d)$.
对下行应用 $\pi_{T,*}$, 并利用 $R^1\pi_{T,*}(\pi_T^*K)=0$, 得到短正合列
\[
0\longrightarrow K=\pi_{T,*}\pi_T^*K
\longrightarrow\pi_{T,*}F_T(d)\longrightarrow\pi_{T,*}Q(d)
\longrightarrow0.
\]
这给出等同 $\pi_{T,*}Q(d)=V$.

若 $Q$ 在 $T$ 上平坦且具有 Hilbert 多项式 $P$, 则
$[F_T\twoheadrightarrow Q]:T\to\Quot^P(F/\mathbb P_S^n)$
给出原映射 $[\pi_{T,*}F_T(d)\twoheadrightarrow V]:T\to\Gr(P(d),\pi_*F(d))$ 的提升, 换言之, $Y\xrightarrow{\sim}T$ 为同构.
上述分析表明, $Y$ 等同于以下 $T$ 的子函子: 
\[
\begin{aligned}
\Sch/T&\longrightarrow\Sets,\\
(T'\to T)&\longmapsto
\begin{cases}
\{*\},&\text{若 $Q_{T'}$ 在 $T'$ 上平坦且 Hilbert 多项式为 $P$; }\\
\varnothing,&\text{否则.}
\end{cases}
\end{aligned}
\]
平坦化分层的存在性\footnote{平坦化分层的存在性可以说是 Quot 概形射影性(\bookref{2.1.3})证明中最深刻的组成部分.}(\bookref{A.2.16})表明, $Y$ 可由 $T$ 的局部闭子概形表出.
\end{proof}
\end{enumerate}
将结论概括如下: 
\begin{proposition}\label{src:2.4.1}
设 $S$ 为 Noether 概形, $\pi:\mathbb P^n\to S$ 为相对射影空间.
若 $F:=\mathcal O_{\mathbb P_S^n}(-l)^{\oplus r}$ 且 $P\in\mathbb Q[z]$, 则存在整数 $m_0$, 使对所有 $d\geq m_0$, 
\[
\Quot^P(F/\mathbb P_S^n)\longrightarrow\Gr(P(d),\pi_*F(d))
\]
为局部闭浸入. 特别地, $\Quot^P(F/\mathbb P_S^n)$ 可由 $S$ 上的拟射影概形表出.
\end{proposition}
\subsection{Hilb 与 Quot 的赋值判别法}\label{src:2.4.2-subsection}
为建立 Quot 的射影性, 只需知道它固有.
\begin{proposition}\label{src:2.4.2}
对每个 Noether 概形 $S$, $\mathbb P_S^n$ 上的凝聚层 $F$ 和多项式 $P\in\mathbb Q[z]$, 函子 $\Quot^P(F/\mathbb P_S^n)$ 满足 $S$ 上的固有性赋值判别法.
即对每个 $S$ 上分式域为 $K$ 的离散赋值环 $R$, $\mathbb P_K^n$ 上 Hilbert 多项式为 $P$ 的每个平坦凝聚商 $F_K\to Q$, 都唯一延拓为 $\mathbb P_R^n$ 上 Hilbert 多项式为 $P$ 的平坦凝聚商 $F_R\to\widetilde Q$.
\end{proposition}
\begin{remark}\label{src:2.4.3}
换言之, 该命题表明每个交换图
\[
\begin{tikzcd}
\operatorname{Spec}K\arrow[r,"Q"]\arrow[d]&
\Quot^P_{X/S}(F)\arrow[d]\\
\operatorname{Spec}R\arrow[r]\arrow[ur,dashed,"\widetilde Q"]&S
\end{tikzcd}
\]
都有唯一提升. 函子和叠的赋值判别法将在 \S\bookref{4.8} 讨论.
\end{remark}
\begin{proof}
记开浸入为 $j:X_K\hookrightarrow X_R$, 将 $\widetilde Q$ 定义为复合
$F_R\to j_*F_K\to j_*Q$ 的像, 其中第一个映射由伴随映射
$F_R\to j_*j^*F_R=j_*F_K$ 给出.
由于 $\widetilde Q$ 为 $j_*Q$ 的子层, 它在 $R$ 上无挠, 故平坦(因为 $R$ 为离散赋值环).
最后, 由于 $Q$ 在 $R$ 上平坦且 $\operatorname{Spec}R$ 连通, 其 Hilbert 多项式为常数.
\end{proof}
\begin{remark}\label{src:2.4.4}
对 $\Hilb^P_{X/S}$, 上述论证变为: 闭子概形 $Z\subseteq X_K$ 的唯一延拓为概形论像
$\widetilde Z=\operatorname{im}(Z\to X_K\hookrightarrow X_R)$.
概形 $\widetilde Z$ 在 $R$ 上平坦, 因为其所有伴随点都位于 $\operatorname{Spec}R$ 的一般点之上.
\end{remark}
\subsection{射影性}\label{src:2.4.3-subsection}
以下结果完成定理~\bookref{2.1.2}~和~\bookref{2.1.3}~的证明.
\begin{theorem}\label{src:2.4.5}
设 $S$ 为 Noether 概形, $\pi:\mathbb P^n\to S$ 为相对射影空间.
若 $F:=\mathcal O_{\mathbb P_S^n}(-l)^{\oplus r}$ 且 $P\in\mathbb Q[z]$, 则存在整数 $m_0$, 使对所有 $d\geq m_0$, 
\[
\Quot^P(F/\mathbb P_S^n)\longrightarrow\Gr(P(d),\pi_*F(d))
\]
为闭浸入. 特别地, $\Quot^P(F/\mathbb P_S^n)$ 可由 $S$ 上的射影概形表出.
\end{theorem}
\begin{proof}
设 $S$ 为 Noether 概形, $\pi:\mathbb P_S^n\to S$ 为相对射影空间, $F$ 为 $\mathbb P_S^n$ 上任意凝聚层.
选择整数 $l,r$ 及满射 $G:=\mathcal O_{\mathbb P_S^n}(-l)^{\oplus r}\twoheadrightarrow F$, 得到函子态射
\[
\begin{aligned}
\Quot^P(F/\mathbb P_S^n)&\longrightarrow\Quot^P(G/\mathbb P_S^n),\\
[F_T\to Q]&\longmapsto[G_T\twoheadrightarrow F_T\twoheadrightarrow Q],
\end{aligned}
\]
上式定义在 $S$-概形 $T$ 上.
由命题~\bookref{2.4.1}, 函子 $\Quot^P(G/\mathbb P_S^n)$ 可由 $S$ 上的拟射影概形表出; 由命题~\bookref{2.4.2}, 它在 $S$ 上固有, 故射影.

为证明 $\Quot^P(F/\mathbb P_S^n)$ 在 $S$ 上射影, 只需证明上述态射可由闭浸入表出.
这归结为如下断言: 对 $S$-概形 $T$ 及 $Q$ 在 $T$ 上平坦的商 $G_T\twoheadrightarrow Q$, 存在闭子概形 $Z\subseteq T$, 使态射 $T'\to T$ 经 $Z$ 分解当且仅当 $G_{T'}\to Q_{T'}$ 经 $F_{T'}$ 分解.
定义 $K=\ker(G_T\to F_T)$, 考虑图
\[
\begin{tikzcd}
0\arrow[r]&K\arrow[r]\arrow[dr]&G_T\arrow[r]\arrow[d,two heads]&F_T\arrow[r]&0\\
&&Q&&
\end{tikzcd}
\]
则 $G_{T'}\to Q_{T'}$ 经 $F_{T'}$ 分解当且仅当 $K_{T'}\to Q_{T'}$ 为零.
由习题~\bookref{1.3.23}, 参数化使 $K_{T'}\to Q_{T'}$ 为零的映射 $T'\to T$ 的 $T$ 的子函子, 可由闭子概形表出.
\end{proof}
取上述证明中的满射 $G=\mathcal O_{\mathbb P_S^n}(-l)^{\oplus r}\twoheadrightarrow F$.
当 $d\gg0$ 时, 有闭浸入的复合
\[
\begin{aligned}
\Quot^P(F/\mathbb P_S^n)&\hookrightarrow\Quot^P(G/\mathbb P_S^n)
\hookrightarrow\Gr(P(d),\pi_*G(d))\\
&\hookrightarrow\mathbb P\left(\bigwedge^{P(d)}\pi_*G(d)\right).
\end{aligned}
\]
在 $\mathbb P_S^n\times_S\Quot^P(F/\mathbb P_S^n)$ 上, 将万有商记为
\[
F_{\Quot^P(F/\mathbb P_S^n)}\longrightarrow Q^{\mathrm{univ}}.
\]
沿上述复合拉回 $\mathcal O(1)$, 得到 $\det p_{2,*}(Q^{\mathrm{univ}}(d))$.
概括如下: 
\begin{corollary}\label{src:2.4.6}
当 $d\gg0$ 时, $\Quot^P(F/\mathbb P_S^n)$ 上的线丛
$\det\bigl(p_{2,*}(Q^{\mathrm{univ}}(d))\bigr)$ 在 $S$ 上相对非常丰沛.
\end{corollary}
\subsection{推广}\label{src:2.4.4-subsection}
我们首先考虑的推广是旗概形.
\begin{exercise}[旗概形]\label{src:2.4.7}
设 $S$ 为 Noether 概形, $F$ 为 $\mathbb P_S^n$ 上的凝聚层, 且可对某些 $l,r$ 写成 $\mathcal O_{\mathbb P_S^n}(-l)^{\oplus r}$ 的商.
对每个多项式序列 $P_1,\ldots,P_k\in\mathbb Q[z]$, 证明函子
\[
\begin{aligned}
\operatorname{Flag}^{P_1,\ldots,P_k}(F/\mathbb P_S^n/S):
\Sch/S&\longrightarrow\Sets,\\
(T\to S)&\longmapsto
\left\{\begin{array}{l}
\text{$\mathbb P_T^n$ 上拟凝聚层的滤过 $0=F_0\subseteq\cdots\subseteq F_k=F_T$, }\\
\text{其中每个商 $F^i/F^{i-1}$ 为有限表现的 $\mathcal O_{\mathbb P_T^n}$-模, }\\
\text{且在 $T$ 上平坦, 具有 Hilbert 多项式 $P_i$}
\end{array}\right\}
\end{aligned}
\]
可由 $S$ 上的射影概形表出.

提示: 将旗概形构造为相对 Quot 概形的迭代.
即考虑 $\operatorname{Flag}:=\operatorname{Flag}^{P_2,\ldots,P_k}(F/\mathbb P_S^n/S)$, 以及
$\mathbb P_S^n\times_S\operatorname{Flag}$ 上的万有滤过
$0=F_{\mathrm{univ},0}\subseteq F_{\mathrm{univ},2}\subseteq\cdots\subseteq p_1^*F$, 
并证明
\[
\operatorname{Flag}^{P_1,\ldots,P_k}(F/\mathbb P_S^n/S)
\cong\Quot^{P_1}(F_{\mathrm{univ},2}/\operatorname{Flag}/S).
\]
\end{exercise}
定理~\bookref{2.1.3}~对 $\mathbb P_S^n$ 上可写成 $\mathcal O_{\mathbb P_S^n}(-l)^{\oplus r}$ 的商的凝聚层 $F$, 证明了 Quot 概形 $\Quot^P(F/\mathbb P_S^n)$ 的存在性与射影性.
现在考虑更一般的情形.
回忆, 态射 $X\to S$ 有三种不同的射影性概念, 按一般性递增列举如下: 
\begin{itemize}
\item $X\to S$ 分解为闭浸入 $X\hookrightarrow\mathbb P_S^n$ 及 $\mathbb P_S^n\to S$; 有时称为 \newterm{$H$-射影}{H-projective}, 因为这是~\cite[II.4]{Har77}~采用的定义; 
\item $X\to S$ 分解为闭浸入 $X\hookrightarrow\mathbb P(E)$ 及 $\mathbb P(E)\to S$, 其中 $E$ 为 $S$ 上的向量丛; 称为\newterm{强射影}{strongly projective}; 
\item $X\to S$ 分解为闭浸入 $X\hookrightarrow\mathbb P(E)$ 及 $\mathbb P(E)\to S$, 其中 $E$ 为 $S$ 上的有限型拟凝聚层(当 $S$ 为 Noether 概形时, 即凝聚层); 这是~\cite[\S II.5]{EGA}~和~\cite[Tag 01W8]{SP}~的定义.
\end{itemize}
若 $S$ 具有\newterm{消解性质}{resolution property}, 即每个凝聚层都是某个向量丛的商, 则射影态射为强射影态射.
每个光滑概形或拟射影概形都具有这一性质.
Hilbert 函子 $\Hilb^P(X/S)$ 和 Quot 函子 $\Quot^P(F/X/S)$ 的定义, 可扩展至任意射影态射 $X\to S$ 及 $X$ 上的有限型拟凝聚层 $F$.
\begin{exercise}\label{src:2.4.8}
设 $X\to S$ 为概形的射影态射, $F$ 为 $X$ 上的凝聚层.
\begin{enumerate}[(a)]
\item (简单)若 $X\to S$ 为 $H$-射影且 $F$ 为 $\mathcal O_X(-l)^{\oplus r}$ 的商, 证明 $\Quot^P(F/X/S)$ 为 $H$-射影.
\item (中等)若 $f:X\to S$ 为强射影且 $F$ 为 $f^*V\otimes\mathcal O_X(-l)$ 的商, 其中 $V$ 为 $S$ 上的向量丛, 证明 $\Quot^P(F/X/S)$ 为强射影.
\item (中等)若 $X\to S$ 射影, 证明 $\Quot^P(F/X/S)$ 为 $S$ 上的固有概形.
\item (困难)若 $X\to S$ 射影且 $F$ 在 $S$ 上平坦, 证明 $\Quot^P(F/X/S)$ 射影.
\item (未解决)若 $X\to S$ 射影, $\Quot^P(F/X/S)$ 是否总是射影的?
\end{enumerate}
另见~\cite[定理 $2.6$]{AK80}~和~\cite[定理 $5.2$]{Nit05}.
\end{exercise}
\begin{exercise}[困难]\label{src:2.4.9}
设 $X\to S$ 为 Noether 概形的强拟射影态射, 即存在局部闭浸入 $X\hookrightarrow\mathbb P(V)$, 其中 $V$ 为 $S$ 上的向量丛.
若 $F$ 为 $X$ 上的凝聚层, $P\in\mathbb Q[z]$ 为多项式, 可以如下修改 Quot 函子: 
\[
\begin{aligned}
\Quot^P(F/X/S):\Sch/S&\longrightarrow\Sets,\\
(T\to S)&\longmapsto
\left\{\begin{array}{l}
\text{$X_T$ 上有限表现的拟凝聚商 $F_T\to Q$, }\\
\text{在 $T$ 上平坦且支集在 $T$ 上固有, }\\
\text{使每个 $t\in T$ 对应的 $X_t$ 上的 $Q_t$}\\
\text{都具有 Hilbert 多项式 $P$}
\end{array}\right\}.
\end{aligned}
\]
类似地, Hilbert 函子 $\Hilb^P(X/S)$ 定义为只参数化支集在基底上固有的闭子概形.
证明 $\Hilb^P_{X/S}$ 和 $\Quot^P_{X/S}(F)$ 可由 $S$ 上的强拟射影概形表出.
另见~\cite[\S$4$]{FGAIV}, \cite[定理 $2.6$]{AK80}~或~\cite[\S$6$]{Nit05}.
\end{exercise}
\begin{remark}\label{src:2.4.10}
若 $X\to S$ 仅为 Noether 概形的分离态射, 则去掉关于 Hilbert 多项式 $P$ 的条件后, 也可如上定义函子 $\Hilb(X/S)$ 和 $\Quot(F/X/S)$.
Artin 代数性公理(\bookref{C.7.4})可用于证明这些函子可由 $S$ 上分离且局部有限型的代数空间表出; 见~\cite[定理 $6.1$]{Art69b}~和~\cite[Tag 09TQ]{SP}, 或作为特殊情形的定理~\bookref{C.7.7}.
Hironaka 的例子给出 $\mathbb C$ 上光滑固有(但非射影)的三维簇 $X$, 使 $\Hilb(X/S)$ 不是概形.
\end{remark}
\subsection{Chow 簇及其他变体}\label{src:2.4.5-subsection}
\newterm{Chow 簇}{Chow variety}是 Hilbert 概形的一种变体, 较易构造, 但函子性质不如后者良好.
Chow 簇 $\operatorname{Chow}_{r,d}(\mathbb P_k^n)$ 参数化 $\mathbb P_k^n$ 上维数为 $r$, 次数为 $d$ 的有效闭链, 可利用有效闭链 $\alpha$ 的 Chow 形式 $\operatorname{Chow}_\alpha$ 构造.
对有效闭链 $\alpha=\sum_i a_i[X_i]$, 可定义
$\operatorname{Chow}_\alpha=\prod_i\operatorname{Chow}_{X_i}^{a_i}$, 
所以只需对维数 $r$, 次数 $d$ 的整闭子概形 $X\subseteq\mathbb P_k^n$ 定义 Chow 形式.
有定理断言, 轨迹
\begin{equation}\label{src:2.4.11}
\{V\in\Gr(r+1,n+1)\mid X\cap V\neq\varnothing\}
\subseteq\Gr(r+1,n+1)
\end{equation}
为次数 $d$ 的不可约超曲面, 且由多项式
\[
\operatorname{Chow}_X\in
H^0(\Gr(r+1,n+1),\mathcal O_{\Gr}(d))
\cong\operatorname{Sym}^d(\bigwedge^{r+1}k^{n+1})
\]
定义, 该多项式在相差非零常数倍的意义下唯一; 见~\cite[\S$3$, 命题 $2.1$--$2.2$]{GKZ08}.

Chow 簇 $\operatorname{Chow}_{r,d}(\mathbb P^n)$ 是 $\mathbb P^n$ 上所有维数 $r$, 次数 $d$ 的有效闭链的 Chow 形式的集合的闭包.
特征 $0$ 中的主要存在性定理断言: 
\begin{itemize}
\item $\operatorname{Chow}_{r,d}(\mathbb P^n)$ 射影且半正规(即每个在剩余域上诱导同构的有限双射态射 $Y\to\operatorname{Chow}_{r,d}(\mathbb P^n)$ 都是同构); 
\item $\operatorname{Chow}_{r,d}(\mathbb P^n)$ 表出半正规 $k$-概形范畴上的一个函子: 
对半正规 $k$-概形 $S$, $\operatorname{Mor}(S,\operatorname{Chow}_{r,d}(\mathbb P^n))$ 等同于 $\mathbb P_S^n$ 上维数 $r$, 次数 $d$ 的良定义的有效代数闭链族; 且
\item 对多项式 $P(t)=d\cdot t^r/r!+\text{(低次项)}$, 存在从 $\Hilb^P(\mathbb P^n)$ 的半正规化出发的 Hilbert--Chow 态射
\[
\Hilb^P(\mathbb P^n)^{\mathrm{sn}}
\longrightarrow\operatorname{Chow}_{r,d}(\mathbb P^n).
\]
它将闭子概形 $Z\subset\mathbb P^n$ 映至闭链 $a_i[Z_i]$, 其中 $Z_i$ 为第 $i$ 个不可约分支的约化概形结构, $a_i$ 为 $Z_i$ 在其一般点处的重数.
\end{itemize}
见~\cite{CW37}, \cite[\S X.6--8]{HP47}~和~\cite{Sam55}, 或现代论述~\cite[\S$5.4$]{GIT}~和~\cite[I.3.21, I.6.1]{Kol96}.
在概述 $\overline M_g$ 的 GIT 构造时, 我们将使用 $\mathbb P^n$ 中曲线的 Chow 簇(即 $r=1$ 的情形), 此时上述断言的证明稍容易些.
\begin{example}[更多变体]\label{src:2.4.12}
还有进一步的变体和推广.
例如, Vistoli 构造了参数化到分离概形 $X$ 的有限非分歧态射的 Hilbert 叠~\cite{Vis91}.
Alexeev 和 Knutson 的分支簇的模参数化从几何约化概形到分离概形 $X$ 的有限态射~\cite{AK10}.
\end{example}
\begin{exercise}[态射概形]\label{src:2.4.13}
对 Noether 概形的强射影态射 $X\to S$ 和 $Y\to S$, 考虑函子
\[
\begin{aligned}
\operatorname{Mor}_S(X,Y):\Sch/S&\longrightarrow\Sets,\\
(T\to S)&\longmapsto\operatorname{Mor}_T(X_T,Y_T),
\end{aligned}
\]
它将 $S$-概形 $T$ 映至 $T$-态射 $X_T\to Y_T$ 的集合.
利用 $\Hilb^P_{X\times_SY/X}$ 中参数化态射 $X\to Y$ 的图像 $X\subseteq X\times_SY$ 的适当闭子概形, 证明 $\operatorname{Mor}_S(X,Y)$ 可由 $S$ 上的射影概形表出.
\end{exercise}
